% Options for packages loaded elsewhere
\PassOptionsToPackage{unicode}{hyperref}
\PassOptionsToPackage{hyphens}{url}
\documentclass[
  12pt,
]{article}
\usepackage{xcolor}
\usepackage{amsmath,amssymb}
\setcounter{secnumdepth}{-\maxdimen} % remove section numbering
\usepackage{iftex}
\ifPDFTeX
  \usepackage[T1]{fontenc}
  \usepackage[utf8]{inputenc}
  \usepackage{textcomp} % provide euro and other symbols
\else % if luatex or xetex
  \usepackage{unicode-math} % this also loads fontspec
  \defaultfontfeatures{Scale=MatchLowercase}
  \defaultfontfeatures[\rmfamily]{Ligatures=TeX,Scale=1}
\fi
\usepackage{lmodern}
\ifPDFTeX\else
  % xetex/luatex font selection
\fi
% Use upquote if available, for straight quotes in verbatim environments
\IfFileExists{upquote.sty}{\usepackage{upquote}}{}
\IfFileExists{microtype.sty}{% use microtype if available
  \usepackage[]{microtype}
  \UseMicrotypeSet[protrusion]{basicmath} % disable protrusion for tt fonts
}{}
\makeatletter
\@ifundefined{KOMAClassName}{% if non-KOMA class
  \IfFileExists{parskip.sty}{%
    \usepackage{parskip}
  }{% else
    \setlength{\parindent}{0pt}
    \setlength{\parskip}{6pt plus 2pt minus 1pt}}
}{% if KOMA class
  \KOMAoptions{parskip=half}}
\makeatother
\setlength{\emergencystretch}{3em} % prevent overfull lines
\providecommand{\tightlist}{%
  \setlength{\itemsep}{0pt}\setlength{\parskip}{0pt}}
\usepackage{bookmark}
\IfFileExists{xurl.sty}{\usepackage{xurl}}{} % add URL line breaks if available
\urlstyle{same}
\hypersetup{
  hidelinks,
  pdfcreator={LaTeX via pandoc}}

\author{SCX}
\date{2026}

\begin{document}

\section{Proposition 2: High-Error Sampling Suboptimality under
Noise}\label{proposition-2-high-error-sampling-suboptimality-under-noise}

\begin{quote}
In the presence of noise, a purely residual-based (max-residual sampling) active sampling strategy is not optimal. High error may originate from learnable difficult states (which should receive more samples) or from intrinsic noise (which should be discarded).
\end{quote}

\begin{center}\rule{0.5\linewidth}{0.5pt}\end{center}

\subsection{1 Statement}\label{statement}

\subsubsection{1.1 Formal Statement}\label{formal-statement}

Let \(\mathcal{D} = \{(x_i, y_i)\}_{i=1}^N\) be the current labeled dataset, with unlabeled pool
\(\mathcal{U}\). For \(x \in \mathcal{U}\), define the residual:

\[r(x) = \ell(f_{\text{current}}(x), \hat{y}(x))\]

where \(\hat{y}(x)\) is the current best prediction.

Define:

\begin{itemize}
\tightlist
\item
  \textbf{Learnable set}
  \(H_{\text{learnable}} = \{x: \mathbb{E}[\ell(f(x), y) \mid x] \text{ can be reduced by more sampling}\}\)
\item
  \textbf{Noise set}
  \(H_{\text{noise}} = \{x: \mathbb{E}[\ell(f(x), y) \mid x] \text{ is irreducible (intrinsic noise)}\}\)
\end{itemize}

\textbf{Max-residual sampling}: \(x^* = \arg\max_{x \in \mathcal{U}} r(x)\)

If \(P(H_{\text{noise}}) > 0\), then the expected budget waste of a pure residual-based sampling strategy
\(\pi_{\text{res}}(x) \propto r_i = \ell(f(x_i), y_i)\)
has a strictly positive lower bound.

\subsubsection{1.2 Assumptions}\label{assumptions}

\begin{enumerate}
\def\labelenumi{\arabic{enumi}.}
\tightlist
\item
  \textbf{Noise existence assumption}: There exists some state \(s_0\) such that
  \(\mathbb{E}[\varepsilon^2 | x \in s_0] \geq \sigma_0^2 > 0\)
\item
  \textbf{Finite sampling budget}: Total labeling budget \(B \ll |\mathcal{U}|\)
\item
  \textbf{Independence across samples}: Label noise is independent across different samples (given \(x\))
\end{enumerate}

\begin{center}\rule{0.5\linewidth}{0.5pt}\end{center}

\subsection{2 Formal Proof}\label{formal-proof}

\subsubsection{2.1
Risk Minimization Framework}\label{risk-minimization-framework}

Define the benefit of sampling strategy \(\pi\) as:

\[J(\pi) = \mathbb{E}_{(x_1^*, \ldots, x_B^*) \sim \pi}\left[\sum_{t=1}^B \Delta R(f_t | x_t^*)\right]\]

where \(\Delta R(f_t | x) = R(f_{t-1}) - R(f_t)\) is the expected risk reduction after labeling \(x\).

Let \(r_i = r(x_i)\), \(\hat{r}_i = \tilde{r}(x_i) - \varepsilon_i\)
be the true residual (unobservable).

\textbf{Lemma 1}: The max-residual strategy \(\pi_{\text{MR}}\) at \(B=1\)
selects:

\[x_{\text{MR}}^* = \arg\max_{i \in \mathcal{U}} \hat{r}_i + \varepsilon_i\]

\textbf{Lemma 2}: The optimal strategy \(\pi^*\) at \(B=1\) selects:

\[x_{\pi^*}^* = \arg\max_{i \in \mathcal{U}} \mathbb{E}[\Delta R(f | x_i)] = \arg\max_{i \in \mathcal{U}} \int \Delta R(f | x_i) dP_{\varepsilon_i}(\varepsilon_i)\]

\textbf{Theorem}: If there exists at least one pair \((i,j)\) such that
\(\hat{r}_i > \hat{r}_j\) but
\(\mathbb{E}[\Delta R(f|x_i)] < \mathbb{E}[\Delta R(f|x_j)]\), then
\(\pi_{\text{MR}} \neq \pi^*\).

\subsubsection{2.2 Proof}\label{proof}

\textbf{Step 1: Relating \(\Delta R\) to \(r\)}

For smooth loss \(L(f(x), y)\), Taylor-expand at \(f_{\text{current}}\):

\[\Delta R(f|x) \approx r(x)^2 \cdot \frac{\partial^2 \mathbb{E}[L]}{\partial f^2}\bigg|_{f_{\text{current}}} + \sigma_y^2(x)\]

where \(\sigma_y^2(x)\) is the label noise variance.

\textbf{Step 2: Noise \(\varepsilon\) introduces bias}

The observed squared residual
\([\tilde{r}(x)]^2 = [r_{\text{true}}(x) + \varepsilon]^2\) satisfies:

\[\mathbb{E}[[\tilde{r}(x)]^2 | x] = r_{\text{true}}(x)^2 + \sigma_\varepsilon^2(x)\]

When \(\sigma_\varepsilon^2(x)\) is large, max-residual
tends to select \textbf{samples with high noise rather than high information content}.

\textbf{Step 3: Construct a counterexample}

Consider two states \(s_1, s_2\), each containing \(n\) points.

State \(s_1\) (learnable): \(y_i = f^*(x_i) + \epsilon_i\), where
\(\epsilon_i \sim N(0, \sigma^2_{\text{low}})\). The current model \(f\) has
systematic bias on \(s_1\):
\(\mathbb{E}[f(x) \mid x \in s_1] = f^*(x) + \mu\), where
\(\mu \gg 0\).

State \(s_2\) (noise): \(y_i = \text{constant} + \eta_i\), where
\(\eta_i \sim N(0, \sigma^2_{\text{high}})\), \(\sigma^2_{\text{high}} \gg \sigma^2_{\text{low}}\).
The current model \(f\) is already optimal on \(s_2\):
\(\mathbb{E}[f(x) \mid x \in s_2] = \text{constant}\).

\textbf{Residual analysis}:

In \(s_1\):
\[r_i^{(1)} = (f(x_i) - y_i)^2 = (\mu - \epsilon_i)^2 = \mu^2 - 2\mu\epsilon_i + \epsilon_i^2\]
\[\mathbb{E}[r^{(1)}] = \mu^2 + \sigma^2_{\text{low}}\]

In \(s_2\):
\[r_i^{(2)} = (f(x_i) - y_i)^2 = \eta_i^2\]
\[\mathbb{E}[r^{(2)}] = \sigma^2_{\text{high}}\]

\textbf{Parameter condition}: Let
\(\mu^2 + \sigma^2_{\text{low}} < \sigma^2_{\text{high}}\). Then the average residual in \(s_2\)
exceeds that in \(s_1\).

\textbf{Max-Residual Sampling behavior}: Select the \(B\) points with the largest global residuals. Since
\(\sigma^2_{\text{high}} \gg \mu^2 + \sigma^2_{\text{low}}\), almost all selected points come from
\(s_2\) (the noise state).

\textbf{Gain analysis}: Sampling additional points from \(s_2\)
does not help reduce prediction error---the model is already optimal there. Each point sampled from \(s_1\)
can reduce bias from \(\mu^2\) to
\(O(\sigma^2_{\text{low}} / n_{s_1})\).

\textbf{Expected gain lower bound}: If Max-Residual samples \(B\) points from \(s_2\):

\[\mathbb{E}[\text{Gain}_{\text{res}}] = O\left(\frac{\sigma^2_{\text{low}}}{n_{s_1}}\right) \cdot \left(\frac{n_{s_1}}{n_{s_1} + n_{s_2}}\right)^B \to 0 \quad \text{as } B \to \infty\]

\textbf{Difference analysis}: Let \(B_{\text{res}}\) be the budget Max-Residual allocates to \(s_1\),
\(B_{\text{state}}\) be the budget State-Level Sampling allocates to \(s_1\). Since Max-Residual is biased toward
\(s_2\), \(B_{\text{res}} \ll B_{\text{state}}\).

Updated expected loss: Max-Residual:
\(\mathbb{E}[\mathcal{L}_{\text{res}}] \approx \frac{\sigma^2_{\text{low}}}{n_{s_1} + B_{\text{res}}} + \sigma^2_{\text{high}}\)
vs. State-Level:
\(\mathbb{E}[\mathcal{L}_{\text{state}}] \approx \frac{\sigma^2_{\text{low}}}{n_{s_1} + B_{\text{state}}} + \sigma^2_{\text{high}}\)

Difference: \(\mathbb{E}[\mathcal{L}_{\text{res}}] - \mathbb{E}[\mathcal{L}_{\text{state}}] \approx \sigma^2_{\text{low}}\left(\frac{1}{n_{s_1} + B_{\text{res}}} - \frac{1}{n_{s_1} + B_{\text{state}}}\right) > 0\).

Therefore
\(\mathbb{E}[\text{Gain}_{\text{state}}] - \mathbb{E}[\text{Gain}_{\text{res}}] > 0\). \(\square\)

\subsubsection{2.3
Supplementary Proof from an Information-Theoretic Perspective}\label{supplementary-proof-information-theoretic}

Let \(I(X; Y | f)\) be the conditional mutual information between unlabeled samples \(X\) and true labels \(Y\)
given the current model. Max-residual sampling maximizes
\(H[\hat{Y} | f, X] = \mathbb{E}[\ell(\hat{Y}, Y)]\), not
\(I(X; Y|f)\).

By the data processing inequality:

\[I(X; Y | f) \leq H[\hat{Y}|f] - H[\hat{Y}|f, X, Y] + H[Y|f] - H[Y|f, X]\]

Under high noise, \(H[Y|f, X]\) increases while \(I(X; Y|f)\) decreases. Max-residual
cannot distinguish ``high uncertainty originating from missing information'' from ``high uncertainty originating from intrinsic noise.''

\begin{center}\rule{0.5\linewidth}{0.5pt}\end{center}

\subsection{3 Noise Score Derivation}\label{noise-score-derivation}

The noise score in SCX is defined as:

\[\text{NoiseScore}(x_i) = r_i \cdot \frac{1}{\rho(s_i) + \varepsilon} \cdot [1 - C(s_i)]\]

where:

\begin{itemize}
\tightlist
\item
  \(r_i = \ell(f(x_i), y_i)\): pointwise residual
\item
  \(\rho(s_i)\): occurrence probability of state \(s_i\) (rare states are downweighted)
\item
  \(C(s_i)\): internal consistency of state \(s_i\)
\end{itemize}

This noise score can be understood as a multi-factor corrected residual:

\begin{enumerate}
\def\labelenumi{\arabic{enumi}.}
\tightlist
\item
  \textbf{Residual factor \(r_i\)}: base error signal
\item
  \textbf{Rarity correction
  \(1/(\rho(s_i)+\varepsilon)\)}: residuals of rare states are inherently unreliable, need downweighting
\item
  \textbf{Consistency correction \([1 - C(s_i)]\)}: when a state is internally consistent, its residuals are more likely to be
  learnable; when inconsistent, more likely to be noise
\end{enumerate}

Relationship between the noise score and Tsybakov's
noise condition: \(C(s) \approx 1 - 2 \cdot \mathbb{E}_{x \in s}[|\eta(x) - 1/2|]\), where
\(\eta(x) = P(Y=1|X=x)\). When \(C(s) \to 1\),
the state interior approximates the Bayes-optimal decision (\(\alpha \to \infty\)); when
\(C(s) \to 0\), labels within the state are nearly random (\(\alpha \to 0\)).

\begin{center}\rule{0.5\linewidth}{0.5pt}\end{center}

\subsection{4 State-Level Alternative}\label{state-level-alternative}

\subsubsection{4.1
State-Level Sampling Criterion}\label{state-level-sampling-criterion}

SCX's state-level active sampling selects:

\[s^* = \arg\max_s \bar{r}(s) \cdot \rho(s) \cdot C(s)\]

where: \(\bar{r}(s)\): average residual of state \(s\) (represents improvement potential),
\(\rho(s)\): occurrence probability of state \(s\) (represents coverage), \(C(s)\): internal consistency of state
\(s\) (represents learnability).

\subsubsection{4.2 Improvement Upper Bound}\label{improvement-upper-bound}

If SCX's state acquisition function \(V(s)\) can correctly distinguish learnable states from noise states (i.e.,
\(L(s)\) is an accurate measure of \(s\)'s learnability), then the expected budget waste of SCX
sampling is at most:

\[\mathbb{E}[\text{Waste}_{\text{SCX}}] \leq B \cdot \epsilon_{\text{est}}\]

where
\(\epsilon_{\text{est}} = P(L(s) \text{ is estimated incorrectly})\). In contrast, Max-Residual's
expected waste lower bound is:

\[\mathbb{E}[\text{Waste}_{\text{res}}] \geq B \cdot P(x \in H_{\text{noise}} \mid \text{highest residual})\]

When
\(P(x \in H_{\text{noise}} \mid \text{highest residual}) \gg \epsilon_{\text{est}}\),
SCX significantly outperforms Max-Residual.

\begin{center}\rule{0.5\linewidth}{0.5pt}\end{center}

\subsection{5 Implications for SCX}\label{implications-for-scx}

\begin{enumerate}
\def\labelenumi{\arabic{enumi}.}
\item
  \textbf{Explains why data poisoning can be detected by SCX}:

  \begin{itemize}
  \tightlist
  \item
    High error + isolated (low \(\rho\)) + state-inconsistent (low \(C\)) \(\to\) noise / poison
  \item
    High error + dense (high \(\rho\)) + state-consistent (high \(C\)) \(\to\) learnable hard
    state
  \end{itemize}
\item
  \textbf{Budget allocation strategy}: Annotation budget should be prioritized for
  states with high \(\bar{r}(s) \cdot \rho(s) \cdot C(s)\), rather than for the samples with the highest residuals.
\item
  \textbf{Relationship to Huber's contamination model}: The noise score \(\text{NoiseScore}(x_i)\)
  serves as the basis for identifying contaminated points in state \(s_i\):
  \[P_{X|s} = (1 - N(s)) \cdot P_{X|s}^0 + N(s) \cdot P_{X|s}^{\text{noise}}\]
  where \(N(s)\) is the noise proportion.
\item
  \textbf{Connection to UCB bandits}: \(V(s)\) can be understood as a generalization of the UCB
  acquisition function for state-arm bandits:
  \[V(s) = \underbrace{\bar{r}(s) \cdot \rho(s)}_{\text{exploitation}} \cdot \underbrace{L(s) \cdot [1-D(s)]}_{\text{exploration}} \cdot \underbrace{\max_m SCX_m(s)}_{\text{expertise}}\]
\end{enumerate}

\begin{center}\rule{0.5\linewidth}{0.5pt}\end{center}

\subsection{References}\label{references}

\begin{enumerate}
\def\labelenumi{\arabic{enumi}.}
\tightlist
\item
  SCX Core Framework Mathematical Analysis. \texttt{01\_SCX\_Core\_Framework\_Math\_Analysis.md} Section
  2.2
\item
  SCX Mathematical Foundations and Proof Construction. \texttt{05\_Math\_Roots\_and\_Proofs.md} Proposition 2
\item
  Tsybakov, A. B. (2004). Optimal aggregation of classifiers in
  statistical learning. \emph{Annals of Statistics}.
\item
  Huber, P. J. (1964). Robust estimation of a location parameter.
  \emph{Annals of Mathematical Statistics}.
\item
  MacKay, D. J. C. (1992). Information-based objective functions for
  active data selection. \emph{Neural Computation}.
\end{enumerate}

\end{document}
